% Results fragment for Section "Geopolitical alignment and the allied coalition"
% (sections/075_alignment_coalition.tex, \label{sec:coalition}, between the opportunity-security
% space and the Strategic Index). It is meant to become most of that section: method (ideal
% points, axis, blocs, inheritance, dyadic alignment, tiers, EU node) and results.
% Generated numbers: paper/results/coalition.json; figures and table: scripts/paper/fig_coalition.R
% (final rerun of 2026-10-12: 57 members = 31 core + 26 partners, 6 EU node-only members,
% 63 economies, 37 network nodes).
% Figures: coalition_axis, coalition_routes, coalition_alignment_matrix, coalition_eu_weights.
% Table: tables/coalition_composition.tex (tab:coalition_composition).
% Citations needed but not in references.bib: USGS Mineral Commodity Summaries 2026, the NEIS Allied
% Industrial Base working paper (October 2026) and the
% OECD-NEA/IAEA Red Book (Uranium 2026); all three are named in the prose without a citation command.

\subsection{Ideal points and the U.S.--China axis}\label{subsec:coalition_axis}

Political alignment is measured from recorded votes in the United Nations General Assembly
(UNGA). We use the ideal points of \citet{bailey2017estimating}, which place every voting state
on one dimension of revealed foreign-policy preference in each session and are estimated with
bridging resolutions that keep positions comparable over time. The headline series is the
all-votes series of the current release \citep[Harvard Dataverse, V39.0, July 2026;][]{voeten2009unga},
estimated from all recorded votes, including votes on separate paragraphs and amendments; the
final-passage series of the same release is used as a check at the end of this section. The
ideal point $\theta_i$ of economy $i$ is its mean over the five latest sessions of the series,
sessions 74 to 78 (2019--2023). The window covers 193 economies that vote in their own right;
Venezuela has ideal points for three of the five sessions and every other economy for all five.

Each economy is placed on the segment between the two anchors,
\begin{equation}\label{eq:coalition_axis}
  s_i \;=\; \min\Bigl\{1,\;\max\Bigl\{0,\;
  \frac{\theta_i-\theta_{\mathrm{CHN}}}{\theta_{\mathrm{USA}}-\theta_{\mathrm{CHN}}}\Bigr\}\Bigr\},
\end{equation}
so that $s_i=0$ is China's position and $s_i=1$ that of the United States
($\theta_{\mathrm{USA}}=2.34$, $\theta_{\mathrm{CHN}}=-1.02$ in the window). Seven economies vote
further from the United States than China does (Cuba, Iran, Nicaragua, North Korea, Syria,
Venezuela and Zimbabwe) and are set to $s_i=0$; none lies beyond the United States. The axis
records where an economy's revealed preferences sit relative to the two anchors. It is not a
measure of alliance, of security commitments or of any bilateral relationship with either anchor.

Economies are grouped into blocs by quartiles of the axis, following \citet{gopinath2025changing}.
The U.S.-leaning bloc is the quarter of voting economies closest to the United States
($s_i\geq 0.511$, the 75th percentile of $s$ over the voting economies other than the United
States), the China-leaning bloc the quarter closest to China ($s_i\leq 0.142$), and the half in
between is nonaligned; the blocs have 49, 95 and 49 members. Membership is hard: the probability
$p_i$ that economy $i$ is on the U.S. side of the divide is 1 in the U.S.-leaning bloc, 0 in the
China-leaning bloc and $\tfrac12$ for a nonaligned economy. The distribution of $s$ is skewed
toward China. The median voting economy sits at $s=0.22$, and 142 of the 193 (73.6 percent) are
closer to China than to the United States. This is a property of the Assembly's voting
arithmetic, not of the economies' economic ties: the United States votes in small minorities on
many recorded resolutions, so its anchor lies at the edge of the distribution. The choice of rule
matters mainly on the China side. A midpoint rule ($p_i=1$ for $s_i>\tfrac12$) would give a
U.S.-leaning group of a similar size, 51 economies, but would classify 142 economies as
China-leaning; the quartile rule keeps the China-leaning bloc to the quarter of the Assembly that
votes most closely with China.

Economies without a vote of their own receive an alignment from a reference map, applied after
the cuts are set so that they do not move the quartiles. Forty-four territories and special
administrative regions inherit their sovereign's bloc (13 from the United Kingdom, 11 from
France, six from the United States, four from the Netherlands, three from Australia, two each
from China and Denmark, and one each from Finland, Norway and New Zealand), and Taiwan, Palestine
and Kosovo are neutral, $p_i=\tfrac12$. All 240 economies of the country universe therefore have
an alignment. Inherited and neutral positions are used only in the bilateral indices; an economy
without its own vote cannot enter the core of the coalition defined below.

For a pair of economies the dyadic alignment is
\begin{equation}\label{eq:coalition_dyadic}
  A_{ij} \;=\; p_i\,p_j + (1-p_i)(1-p_j),
\end{equation}
the probability that $i$ and $j$ fall on the same side of the U.S.--China divide if each is on
the U.S. side independently with probability $p$. With hard blocs $A_{ij}$ takes three values:
1 for two economies in the same bloc, 0 for economies in opposite blocs and $\tfrac12$ whenever
either economy is nonaligned. Two nonaligned economies also have $A_{ij}=\tfrac12$: the measure
concerns the U.S.--China cleavage only, not the similarity of two voting records in general.
Alignment enters the rest of the analysis in two places. $A_{ij}$ is one of the seven components
of the Friendshore index, with weight 2 of a total component weight of 11, that is 18.2 percent
(Section~\ref{subsec:friendshore}); through Friendshore it also enters the edge term of the
network design (Section~\ref{sec:network}). And the blocs define the coalition $\Coalition$.
Figure~\ref{fig:coalition_axis} shows the axis for every voting economy, by coalition tier.

\begin{figure}[tbp]
  \centering
  \includegraphics[width=\textwidth]{coalition_axis}
  \caption{Positions on the U.S.--China axis and coalition tiers. Each point is one of the 193
  economies with their own votes in the United Nations General Assembly, at its position
  $s_i$ of equation~\eqref{eq:coalition_axis}: the mean all-votes ideal point of
  \citet{bailey2017estimating} over sessions 74--78 (2019--2023), projected on the segment from
  China ($s=0$) to the United States ($s=1$) and clamped to $[0,1]$. Vertical lines are the
  quartile cuts that define the blocs ($s\leq 0.142$ China-leaning, $s\geq 0.511$ U.S.-leaning;
  49, 95 and 49 economies). Rows are coalition tiers: the 31 core members, the 26 partners, the
  six EU members that belong to the network only through the EU node, and the 130 other voting
  economies; within a row, points are stacked to avoid overlap and the vertical position carries
  no information. Filled markers are major suppliers (at least 5 percent of world supply in one
  measure of an Upstream or Midstream stage of a headline technology, or at least 10 percent of
  world reserves of a headline mineral, fuel or uranium), whatever their bloc. G20 economies and
  selected other cases are labelled. Territories, which inherit their sovereign's bloc, and
  Taiwan, Palestine and Kosovo, which are neutral, have no vote of their own and are not shown.
  Source: \texttt{geo\_ideal\_points\_tbl.rds}, \texttt{coalition\_tbl.rds},
  \texttt{major\_supplier\_shares\_tbl.rds}.}
  \label{fig:coalition_axis}
\end{figure}

\subsection{A coalition in two tiers}\label{subsec:coalition_tiers}

The coalition $\Coalition$ is derived by rule rather than chosen. Its core is the U.S.-leaning
bloc restricted to economies with their own UN votes and a latest GDP of at least USD 100 billion
\citep[current dollars, 2025 for every member;][]{worldbank2026wdi}. It has 31 members, and every
U.S.-leaning economy above the GDP floor belongs to it. The 18 U.S.-leaning economies below the
floor include six EU members (Cyprus, Estonia, Latvia, Lithuania, Malta and Slovenia) and small
European and Pacific states. Puerto Rico inherits the U.S. position and its GDP exceeds the floor,
but it has no vote of its own and is not a member.

An economy outside the core is a partner if it qualifies by either of two routes. The agreement
route admits members of a U.S. free trade agreement, the Indo-Pacific Economic Framework (IPEF),
the Minerals Security Partnership (MSP), the G7 or NATO whose GDP is at least USD 100 billion,
whatever their bloc. The supplier route admits, without a GDP floor, any economy that is not
China-leaning and is a major supplier of an Upstream or Midstream stage of a headline technology:
it holds at least 5 percent of world supply in one measure of the stage, or at least 10 percent
of world reserves of a resource the stage extracts. The supply measures are exports of the
stage's HS6 basket \citep[2023;][]{growthlab2019trade}, oil, gas and coal production
\citep[2025;][]{ei2026statistical}, uranium mine production (2024), mine production and refining
of the minerals a technology uses \citep[2025;][with manganese, zinc and platinum-group metals
from the Energy Institute]{iea2026cmdata} and clean-technology manufacturing capacity (2024). The
reserves are the critical-mineral reserves of the U.S. Geological Survey's \emph{Mineral
Commodity Summaries 2026} at the end of 2025, proved oil, gas and coal reserves at the end of 2020
\citep{ei2026statistical} and identified uranium resources recoverable below USD 260/kgU from the
OECD-NEA/IAEA \emph{Uranium 2026} (the Red Book). The reserves bar is twice the supply bar because
a stock is not yet supply. The route uses the same reserves data as the Reserves category of the
Energy Security pillar (Section~\ref{subsec:es_reserves}), with one consequence worth stating:
holders whose reserves USGS does not quantify (entries NA, W or s) count as zero, so their shares,
and the coalition's, are measured on the quantified holdings only.

Twenty-six economies are partners: ten qualify through agreements only, nine as suppliers only and
seven by both routes (Table~\ref{tab:coalition_composition}, Figure~\ref{fig:coalition_routes}).
Twenty-three partners are nonaligned and three are China-leaning: Indonesia and Viet Nam, members
of IPEF, and Oman, a U.S. free-trade partner. The agreement route has no bloc condition; Indonesia
and Viet Nam also pass the supplier test, which their bloc would not let them use. No partner is
U.S.-leaning, since every U.S.-leaning economy above the floor is already in the core. NATO is the
most frequent qualifying agreement among the members (24), followed by the U.S. free trade
agreements (14), the MSP (13), IPEF (12) and the G7 (7). Twenty-eight of the 31 core members
belong to at least one of them (Austria, Ireland and Ukraine to none), and 13 core members also
pass the supplier test. Sixteen members of the five agreements fall below the GDP floor and are
not partners; the largest are Lithuania (USD 95 billion), which enters the network through the EU
node, and Panama (USD 91 billion). Together the 57 members account for 73.4 percent of world GDP
(57.7 percent the core and 15.7 percent the partners), and 73.7 percent with the six EU members
that enter only through the EU node; China accounts for 16.6 percent.

Forty economies pass the supplier test: 13 U.S.-leaning, 16 nonaligned and 11 China-leaning.
Every one of them that is not China-leaning is a member. The nine China-leaning major suppliers
are kept out: China, Russia, Iran, Kazakhstan, Algeria, Uzbekistan, Venezuela, Myanmar and
Zimbabwe. Their largest qualifying shares range from 94.4 percent of world graphite refining
(China) and 45.1 percent of the exports of the Nuclear Upstream basket of uranium ores and
concentrates (Kazakhstan) to 5.1 percent of the exports of the Gas Upstream basket (Algeria);
Myanmar qualifies through 31.0 percent of the exports of the Wind Upstream basket of rare-earth
compounds, and Venezuela through its oil reserves alone, 17.5 percent of the world's proved
reserves. Ten members hold at least 10 percent of world reserves of some resource (the United
States, Australia, India, Brazil, Saudi Arabia, Argentina, South Africa, Chile, Qatar and the
Democratic Republic of the Congo), but for none of them is the reserves leg decisive: each also
passes a supply measure. Brazil, the case that motivated the leg, holds 23.9 percent of world
graphite reserves, 16.7 percent of manganese, 14.6 percent of rare earths and 11.4 percent of
nickel, and since silicon metal entered the Solar Upstream basket it also supplies 18.4 percent
of world exports of that basket. Among all economies the reserves leg is decisive only for
Venezuela, which is kept out. The bar binds for a small group of holders. Six economies that are
not China-leaning hold between 5 and 10 percent of a world reserve without qualifying otherwise:
Madagascar (8.7 percent of graphite), Mozambique (8.1 percent of graphite), Kuwait (5.9 percent of
oil), Tanzania (5.8 percent of graphite), Niger (5.7 percent of uranium) and New Caledonia (5.1
percent of nickel); Iraq, Turkmenistan and the United Arab Emirates fall in the same range but are
China-leaning. Admitting the Forum of the MSP as a further agreement, a sensitivity, adds
Kazakhstan and Uzbekistan, both China-leaning, and Ecuador, which is nonaligned.

% Generated by scripts/paper/fig_coalition.R from data/processed/coalition_tbl.rds,
% geo_ideal_points_tbl.rds and major_supplier_shares_tbl.rds
\begin{table}[htbp]
\centering
\begin{threeparttable}
\caption{Composition of the allied coalition}
\label{tab:coalition_composition}
\footnotesize
\begin{tabularx}{\textwidth}{@{}>{\raggedright\arraybackslash}p{0.27\textwidth} r >{\raggedright\arraybackslash}X@{}}
\toprule
Group & $n$ & Economies \\
\midrule
Core: U.S.-leaning bloc, GDP $\geq$ USD 100bn & 31 & United States, Germany, Japan, United Kingdom, France, Italy, Canada, Spain, Korea, Australia, Netherlands, Poland, Belgium, Ireland, Sweden, Israel, Austria, Norway, Denmark, Romania, Czechia, Portugal, Finland, Greece, New Zealand, Hungary, Ukraine, Slovakia, Bulgaria, Croatia, Luxembourg \\
Partners, agreement route only & 10 & T\"urkiye (NATO), Indonesia (IPEF), Singapore (IPEF, U.S. FTA), Thailand (IPEF), Viet Nam (IPEF), Morocco (U.S. FTA), Dominican Rep. (U.S. FTA), Guatemala (U.S. FTA), Oman (U.S. FTA), Costa Rica (U.S. FTA) \\
Partners, supplier route only & 9 & Brazil (Solar Up 18\% [+3]), Saudi Arabia (Oil Up 13\%), Argentina (Batteries Up 6\% [+1]), South Africa (Hydrogen Up 52\% [+3]), Qatar (Gas Up 14\%), Ghana (Batteries Up 10\% [+2]), DR Congo (Batteries Up 66\% [+9]), Gabon (Batteries Up 25\% [+2]), Namibia (Nuclear Up 12\%) \\
Partners, both routes & 7 & India (IPEF, MSP), Mexico (U.S. FTA), Philippines (IPEF), Malaysia (IPEF), Colombia (U.S. FTA), Chile (U.S. FTA), Peru (U.S. FTA) \\
EU members in the EU node only & 6 & Lithuania, Slovenia, Latvia, Estonia, Cyprus, Malta \\
Major suppliers kept out (China-leaning) & 9 & China (Batteries Mid 94\% [+16]), Russia (Hydrogen Up 29\% [+6]), Iran (Gas Up 6\% [+1]), Kazakhstan (Nuclear Up 45\%), Algeria (Gas Up 5\%), Uzbekistan (Nuclear Up 8\%), Venezuela (Oil reserves 18\%), Myanmar (Wind Up 31\% [+1]), Zimbabwe (Batteries Up 9\% [+2]) \\
Reserves holders between 5\% and 10\%, not members & 6 & Madagascar (Graphite 8.7\%), Mozambique (Graphite 8.1\%), Kuwait (Oil 5.9\%), Tanzania (Graphite 5.8\%), Niger (Uranium 5.7\%), New Caledonia (Nickel 5.1\%) \\
Added by the MSP Forum sensitivity & 3 & Kazakhstan, Uzbekistan, Ecuador \\
Hand-picked list of the previous draft: not members & 2 & UAE, Kenya \\
Hand-picked list of the previous draft: members & 21 & 11 in the core, 10 partners \\
\bottomrule
\end{tabularx}
\begin{tablenotes}[flushleft]\footnotesize
\item \emph{Notes:} Economies are ordered by GDP within a group. Core: the U.S.-leaning quartile bloc of the all-votes ideal points of \citet{bailey2017estimating} over UNGA sessions 74--78 (2019--2023), $s\geq 0.511$, with the economy's own UN votes and latest World Bank GDP of at least USD 100bn. Partners (not core): a member of a U.S. free trade agreement, IPEF, the Minerals Security Partnership (MSP), the G7 or NATO with GDP of at least USD 100bn (qualifying agreements in parentheses), or, without a GDP floor, an economy that is not China-leaning and holds at least 5\% of world supply in one measure of an Upstream or Midstream stage of a headline technology (exports of the stage's HS6 basket, fuel or uranium production, mineral mining or refining, or manufacturing capacity) or at least 10\% of world reserves of a headline mineral, fuel or uranium. For the supplier route, the parentheses give the largest qualifying share: the technology--stage cell of a supply measure, or the resource when only reserves qualify, and in brackets the number of further qualifying cells. The six EU members below the GDP floor belong to the coalition only through the EU. The MSP Forum sensitivity admits Forum members by the agreement route. The previous draft's 23 members were fixed by hand.
\end{tablenotes}
\end{threeparttable}
\end{table}


\begin{figure}[tbp]
  \centering
  \includegraphics[width=\textwidth]{coalition_routes}
  \caption{Routes into the coalition. Rows are the 31 core members (left) and the 26 partners
  (right), ordered by GDP. Gray tiles mark membership of the five agreements of the agreement
  route (a U.S. free trade agreement, IPEF, the Minerals Security Partnership, the G7 and NATO).
  Coloured tiles mark the supplier test: at least 5 percent of world supply in one measure of an
  Upstream (blue) or Midstream (orange) stage of a headline technology, or at least 10 percent of
  world reserves of a headline mineral, fuel or uranium (light blue). Outlined tiles mark
  economies that pass the test but are China-leaning ($\dagger$) and so cannot use the supplier
  route; Indonesia and Viet Nam enter through IPEF and Oman through its U.S. free trade agreement.
  The core requires neither route: its tiles show which core members would also qualify as
  partners. The text to the right gives the largest qualifying share, as the technology--stage
  cell of a supply measure or, where only reserves qualify, the resource, with the number of
  further qualifying technology--stage cells in parentheses. Supply measures: Atlas export shares
  of the stage's HS6 basket (2023), EI production (2025), uranium mine production (2024), IEA and
  EI mineral mining and refining (2025), IEA manufacturing capacity (2024); reserves: USGS (end
  of 2025), EI (end of 2020), identified uranium resources below USD 260/kgU. Source:
  \texttt{coalition\_tbl.rds}, \texttt{major\_supplier\_shares\_tbl.rds},
  \texttt{economic\_security\_agreements.csv}.}
  \label{fig:coalition_routes}
\end{figure}

Two economies that a reader might expect to find in the coalition are not members. Switzerland
is nonaligned, at $s=0.508$, three thousandths below the U.S.-leaning cut; it belongs to none of
the five agreements, and its largest share of world supply in any Upstream or Midstream measure
is 1.1 percent. Before the review of the HS6 baskets it qualified as a supplier through its
exports of electrical energy, which are no longer part of the Electric Grid basket
(Section~\ref{subsec:data_crosswalk}). Taiwan has no UN vote and is therefore neutral; it appears
in none of the agreement lists, has no World Bank GDP, and its largest share in any measure is
1.5 percent, of the exports of the Solar Midstream basket. Its strength in semiconductors lies
outside the energy baskets. Both remain in the country universe and in the bilateral indices,
where they are rated as partners of the coalition, but neither is a node of the network.

The coalition's command of supply differs sharply between extraction and processing. Across the
18 Upstream and Midstream stages that have an HS6 basket, the 63 economies of the network account
for a median 70.7 percent of world exports, from 45.2 percent in Wind Upstream, where Myanmar's
rare-earth exports weigh, to 88.4 percent in Batteries Upstream, and China for a median 17.6
percent, with a maximum of 40.9 percent in Batteries Midstream. The coalition's share exceeds
China's in every one of the 18 stages, but in two Upstream stages the excluded China-leaning
suppliers other than China hold large shares: 31.3 percent of Wind Upstream exports, against the
coalition's 45.2 percent, and 47.6 percent of Nuclear Upstream exports, against 50.0 percent. Export shares
flatter the coalition, because they count trade among its members, intra-EU trade included, and
re-exports. Mineral production data show the division more clearly. Coalition economies mine 81.9 percent of the world's cobalt, 79.4 percent
of its nickel, 60.0 percent of its lithium and 54.5 percent of its copper, but 3.6 percent of its
graphite and 14.3 percent of its rare earths; they refine 17.4 percent of cobalt, 23.2 percent of
copper, 24.6 percent of lithium, 5.5 percent of graphite and 12.6 percent of rare earths, against
China's 75.6, 46.8, 73.7, 94.4 and 83.9 percent \citep{iea2026cmdata}. Nickel is the exception,
with 55.7 percent of refining in the coalition. Reserves follow extraction: coalition economies
hold 75.8 percent of quantified world cobalt reserves, 77.3 percent of lithium, 85.5 percent of
manganese, 78.8 percent of nickel and 83.9 percent of platinum-group metals, but 33.3 percent of
rare earths (China 58.4 percent) and 35.7 percent of graphite (China 32.3 percent). In oil and gas
the excluded China-leaning suppliers other than China hold almost as much as the coalition or more
(35.3 against 35.7 percent of oil reserves, 43.4 against 30.8 percent of gas). These are
descriptive shares of world totals, not the network design's inputs; they indicate where an
allied industrial base built from this coalition is long, in mining, and where it is short, in
mineral processing.

\subsection{Network nodes and the European Union}\label{subsec:coalition_nodes}

The EU-27 enters the network design as a single node, because the Union conducts a common
commercial policy and its members' supply chains are integrated in a single market. Twenty-one EU
members are core members in their own right; the six below the GDP floor belong to the network
only through the EU node. They account for 1.6 percent of EU GDP; the largest weight any of them
carries in a stage is Estonia's 25.5 percent of the EU's Wind Upstream exports. The network therefore has 37
nodes representing 63 economies: 11 core nodes (the United States, the EU, Japan, the United
Kingdom, Canada, Korea, Australia, Israel, Norway, New Zealand and Ukraine) and 26 partner nodes.
The EU node's GDP, USD 21.2 trillion, is the second largest after that of the United States.
How the node's demand enters the design, where it is based on ranks and is therefore compressed
relative to the EU's share of coalition GDP, is set out in Section~\ref{subsec:network_demand}.

The members' inputs are aggregated into the node with stage-specific weights
(Figure~\ref{fig:coalition_eu_weights}). For each of the 22 Upstream and Midstream cells, the
weights are the members' shares in the first usable basis: production (Energy Institute) for
Coal, Gas and Oil Upstream (3 cells), exports of the stage's HS6 basket for 15 cells, and GDP for
the remaining 4 (Geothermal and Green Hydrogen, Upstream and Midstream), the stages without a
basket. A basis is usable only if it lists the members individually and covers at least half of
the EU total, and only if the EU holds at least 0.1 percent of the world total; the IEA
manufacturing and mineral data, which report the EU as a bloc or list only leading producers, do
not qualify. The GDP basis is therefore a fallback that carries no information about the stage.
Weights are concentrated. Germany carries the largest weight in 11 of the 22 cells (23.8 percent
in the GDP-weighted cells, 38.9 percent of the Electric Vehicles Midstream basket exports), France
and Poland in four each, and Finland (Electric Vehicles Upstream, 27.4 percent), Italy (Oil
Upstream, 38.5 percent) and Romania (Gas Upstream, 29.5 percent) in one each. The new Upstream
baskets concentrate the weights most: France, which mines no uranium, accounts for 97.9 percent
of the EU's exports of uranium ores and concentrates (Nuclear Upstream), trade that reflects its
conversion and trading activity rather than mine output, and for 46.5 percent of silicon-metal
exports (Solar Upstream);
Estonia, a member below the GDP floor, carries 25.5 percent of the Wind Upstream basket of
rare-earth compounds. The median Herfindahl index of the weights is 0.34 under the production
basis, 0.139 under the export basis and 0.117 under the GDP basis, with a range from 0.081 to
0.959 across cells. The production weights rest on small totals: the EU produces 1.6, 0.8 and 0.3
percent of world coal, gas and oil, and the listed members cover 95, 89 and 64 percent of the
EU's own totals.

\begin{figure}[tbp]
  \centering
  \includegraphics[width=\textwidth]{coalition_eu_weights}
  \caption{Member weights within the EU network node. Columns are the 27 EU members, ordered by
  GDP; an asterisk marks the six members below the GDP floor, which enter the network only through
  the node. Rows are the 22 Upstream and Midstream technology--stage cells, grouped by the basis
  of the weights: members' production (Energy Institute; 3 cells), members' shares of world
  exports of the stage's HS6 basket (Atlas, 2023; 15 cells) or, for the four stages without a
  basket, GDP shares (Geothermal and Green Hydrogen, Upstream and Midstream; Germany 0.238, France 0.158, Italy
  0.120, Spain 0.090, the Netherlands 0.063). Weights in a row sum to one; weights of at least
  0.15 under the production and export bases are printed. Source:
  \texttt{eu\_stage\_weights\_tbl.rds}.}
  \label{fig:coalition_eu_weights}
\end{figure}

\subsection{Alignment within the coalition}\label{subsec:coalition_dyadic}

Because the blocs are hard, alignment within the coalition takes a simple form
(Figure~\ref{fig:coalition_alignment_matrix}). Of the 666 pairs of network nodes, 8.7 percent
have $A_{ij}=1$: the 55 pairs of core nodes and the three pairs among the China-leaning partners.
A further 86.3 percent have $A_{ij}=\tfrac12$, because one of the two is a nonaligned partner,
and 5.0 percent have $A_{ij}=0$: each of the 11 core nodes paired with Indonesia, Viet Nam or
Oman. Counting the EU members individually, the 1,596 pairs of members split 29.3, 64.8 and 5.8
percent. Every pair of core members has $A_{ij}=1$, every nonaligned partner has $A_{ij}=\tfrac12$
with every core member, and the three China-leaning partners have $A_{ij}=0$ with the core and
$A_{ij}=1$ with China and Russia. Within the coalition the alignment component of Friendshore
therefore separates core suppliers from partner suppliers and penalises the three China-leaning
partners as suppliers to the core, but it does not distinguish among the 23 nonaligned partners.
Variation among them in Friendshore comes from the other six components.

\begin{figure}[tbp]
  \centering
  \includegraphics[width=\textwidth]{coalition_alignment_matrix}
  \caption{Dyadic alignment among the network nodes. Each cell is $A_{ij}=p_ip_j+(1-p_i)(1-p_j)$
  of equation~\eqref{eq:coalition_dyadic} for the row and column economies, where $p_i$ is 1 in
  the U.S.-leaning bloc, $\tfrac12$ for a nonaligned economy and 0 in the China-leaning bloc; with
  hard blocs $A_{ij}$ is 1 for the same bloc, 0 for opposite blocs and $\tfrac12$ when either
  economy is nonaligned. The 37 network nodes, 11 core and 26 partner, are followed by Russia and
  China for reference; nodes are ordered by tier and, within a tier, by $s$ (for the EU node, the
  GDP-weighted mean of its members, 0.677; all 27 members are U.S.-leaning, so the node's
  alignment with any other economy equals that of each of its members). Source:
  \texttt{geo\_alignment\_tbl.rds}, \texttt{allied\_network\_nodes.csv}.}
  \label{fig:coalition_alignment_matrix}
\end{figure}

\subsection{Sensitivity of the alignment and comparison with the hand-picked list}\label{subsec:coalition_sensitivity}

Two alternative series test the axis. Using only the latest session of the all-votes series
(session 78, 2023), the Spearman rank correlation of $s$ with the headline five-session mean is
0.966 over 192 economies, and 89.1 percent of economies stay in the same bloc. The coalition
keeps 57 members but changes at the margin: New Zealand falls just below the U.S.-leaning cut
($s=0.583$ against 0.584) and becomes a partner through IPEF, Saudi Arabia becomes China-leaning
and loses its supplier route, and Myanmar becomes nonaligned and enters as a supplier. India also
becomes China-leaning but remains a partner through IPEF and the MSP. Using the final-passage
series over its latest five sessions (76--80, 2021--2025), the Spearman rank correlation is 0.965
over 193 economies and 89.6 percent stay in the same bloc. The core is unchanged. The cuts,
however, move to 0.446 and 0.055: China-leaning economies cluster more tightly around China in
final-passage votes, and Russia ($s=0.117$), Kazakhstan (0.056) and Myanmar (0.163) leave the
China-leaning quartile and would enter the coalition as major suppliers, for 60 members; Malaysia
becomes China-leaning but stays through IPEF. The core, defined by the U.S.-leaning quartile and
the GDP floor, is robust to the choice of series. The exclusion of China-leaning suppliers is the
fragile element of the rule, and Russia's exclusion in particular depends on using all recorded
votes rather than final passage only.

The derived coalition replaces the 23 members fixed by hand in the previous draft. Twenty-one of
them remain members, 11 in the core (the United States, Canada, Japan, Australia, Korea, the
United Kingdom, Germany, France, Italy, Norway and Denmark) and 10 as partners (India, Mexico,
Brazil, Saudi Arabia, South Africa, Indonesia, Viet Nam, Argentina, Morocco and Chile). Two leave:
the United Arab Emirates and Kenya, both China-leaning ($s=0.132$ and 0.142), in none of the five
agreements and without a qualifying share (the Emirates' largest supply share is 4.2 percent and
its oil reserves are 5.6 percent of the world's; Kenya has none). Brazil, which is nonaligned and
belongs to none of the five agreements, remains a member through the supplier route, on both its
reserves and its silicon-metal exports. Thirty-six economies join: 17 EU members, which with Denmark, France,
Germany and Italy form the EU node; Israel, New Zealand and Ukraine in the core; and 16 partners
(Colombia, Costa Rica, the Democratic Republic of the Congo, the Dominican Republic, Gabon, Ghana,
Guatemala, Malaysia, Namibia, Oman, Peru, the Philippines, Qatar, Singapore, Thailand and
T\"urkiye). Every member has a Policy index, so every member can carry a Strategic Index pick
(Section~\ref{sec:strategic}): 215 economies are eligible, 65 on both the IEA policy database and
the Climate Action Tracker rating and 150 on the IEA database alone, and 11 members are scored on
the IEA database alone (Israel, Ukraine, the Democratic Republic of the Congo, the Dominican
Republic, Gabon, Ghana, Guatemala, Malaysia, Namibia, Oman and Qatar). Under a rule that required
both components these 11 members would have no Policy index and no pick.

\paragraph{The rule and a functional definition of the coalition.} The working paper on the Allied
Industrial Base of the New Energy Industrial Strategy programme (October 2026) argues that
``allied'' is better treated as a functional than as a treaty category: partners would be assessed
on reliability and contractual performance, economic complementarity, strategic alignment,
willingness to share burdens, technology-security requirements and contribution to the resilience
of the system, and membership could vary from one sector compact to another. The derived coalition
implements part of that programme. The axis measures strategic alignment, the supplier route
measures a partner's contribution to world supply and reserves in the stages the network designs,
and the agreement route admits economies whose economic-security ties with the United States are
already institutionalised, whatever their bloc. The rule does not measure reliability, the rule of
law, burden sharing or technology-security requirements, for which no comparable cross-country
measures enter the pipeline, and it defines one coalition for all 22 Upstream and Midstream stages
rather than one for each. The supplier route nonetheless gives the boundary a sectoral margin, since
nine partners qualify only as suppliers of particular stages (Figure~\ref{fig:coalition_routes}),
and the sensitivity analyses above show that the boundary moves mainly among the China-leaning
suppliers. A coalition defined stage by stage, which could admit a resource holder to a minerals
compact without admitting it to a compact on sensitive grid equipment, is a natural extension that
the network design could accommodate, and we leave it to future work.
